% ═══════════════════════════════════════════════════════════════════════
%  interpretability_meets_safety.tex
%  "Use of AI Interpretability for AI Safety" --- 30-min presentation
%  Based on: Lee et al. (EMNLP 2025) --- Interpretation Meets Safety
%  Styled to match safety.tex UCR theme
% ═══════════════════════════════════════════════════════════════════════

% ─── TITLE FRAME ─────────────────────────────────────────────────────────────
\begin{frame}[plain]
  \begin{tikzpicture}[remember picture, overlay]
    % Navy header band
    \fill[UCRnavy] (current page.north west)
      rectangle ([yshift=-0.38\paperheight]current page.north east);
    % Gold accent
    \fill[UCRnavy] ([yshift=-0.38\paperheight]current page.north west)
      rectangle ([yshift=-0.41\paperheight]current page.north east);
    % Title text
    \node[anchor=center, text=white, font=\LARGE\bfseries,
          text width=0.88\paperwidth, align=center]
      at ([yshift=0.24\paperheight]current page.center)
      {Interpretation Meets Safety};
    \node[anchor=center, text=white!80!UCRnavy, font=\large,
          text width=0.78\paperwidth, align=center]
      at ([yshift=0.13\paperheight]current page.center)
      {How Interpretability Methods Understand and Enhance LLM Safety};
  \end{tikzpicture}

  \vspace{5.2em}
  \begin{columns}[T]
    \column{0.54\textwidth}
    {\footnotesize\bfseries\navy{Based on:}}\\[3pt]
    {\scriptsize Seongmin Lee, Aeree Cho, Grace C.\ Kim, ShengYun Peng,\\
    Mansi Phute, Duen Horng Chau (Georgia Tech)\\[4pt]
    \textit{``Interpretation Meets Safety: A Survey on Interpretation Methods\\
    and Tools for Improving LLM Safety''}\\[3pt]
    \textbf{EMNLP 2025} --- Nearly 70 works surveyed}

    \column{0.44\textwidth}
    \begin{ucrinfo}
      \scriptsize\bfseries My Role in Today's Presentation:\\[3pt]
      \normalfont\scriptsize
      My teammates covered \textbf{what} interpretability is, and \textbf{what} safety risks exist. I connect the two: how interpretability \emph{actually} helps us build safer LLMs.
    \end{ucrinfo}
  \end{columns}
\end{frame}


% ═══════════════════════════════════════════════════════════════════
%  SLIDE 1 --- WHY BRIDGE INTERPRETABILITY AND SAFETY?
% ═══════════════════════════════════════════════════════════════════
\begin{frame}{Why Bridge Interpretability and Safety?}
  \begin{columns}[T]
    \column{0.48\textwidth}
    {\small\bfseries\navy{The Problem with Existing Safety Methods}}
    \begin{itemize}\footnotesize\setlength\itemsep{4pt}
      \item RLHF, DPO, Constitutional AI work empirically ---
        \textbf{but without understanding why}
      \item Models can be ``deceptively aligned'': pass safety tests
        without internalizing safety (Hubinger et al.)
      \item Fine-tuning on just 100 harmful examples can erase
        all alignment
      \item We are patching symptoms, not fixing root causes
    \end{itemize}
    \vspace{0.4em}
    \begin{ucrwarn}
      \scriptsize\bfseries The Gap in Prior Surveys:\\[3pt]
      \normalfont\scriptsize
      Surveys focus on \emph{either} interpretability \emph{or} safety --- not how the former informs the latter. This is the first survey to bridge both systematically.
    \end{ucrwarn}

    \column{0.48\textwidth}
    {\small\bfseries\tealc{The Insight: Mechanistic Understanding = Better Fixes}}
    \begin{ucrcard}[UCRnavy]{Surgery, Not Chemotherapy}
      \scriptsize
      If we can \textbf{locate} the mechanism behind harmful behavior inside the model, we can \textbf{fix it precisely}:\\[4pt]
      $\bullet$\ Find the ``sycophancy'' direction $\to$ steer it away\\
      $\bullet$\ Locate jailbreak-enabling attention heads $\to$ prune them\\
      $\bullet$\ Trace hallucination to training data $\to$ remove it\\
      $\bullet$\ Find bias SAE neurons $\to$ clamp them to zero
    \end{ucrcard}
    \vspace{0.3em}
    \begin{ucrdark}[UCRnavy]
      \centering\scriptsize\bfseries
      Survey covers nearly 70 works across four safety concerns:\\[2pt]
      \normalfont\scriptsize
      Hallucination · Jailbreaks \& Harm · Bias · Privacy Leakage
    \end{ucrdark}
  \end{columns}
\end{frame}


% ═══════════════════════════════════════════════════════════════════
%  SLIDE 2 --- UNIFIED FRAMEWORK OVERVIEW
% ═══════════════════════════════════════════════════════════════════
\begin{frame}[fragile]{Unified Framework: Interpret $\to$ Enhance $\to$ Tools}
  \vspace*{0.25cm}
  \begin{tikzpicture}[
    box/.style={rectangle, draw=#1, fill=#1!12, text=#1,
                font=\small\bfseries, minimum width=3.2cm,
                minimum height=0.9cm, align=center, rounded corners=2pt},
    arr/.style={-{Stealth[length=5pt]}, ultra thick, #1},
    lbl/.style={font=\tiny\itshape, text=UCRmgray, align=center, text width=3.2cm},
    wfbox/.style={rectangle, draw=UCRmgray, fill=UCRlgray, rounded corners=1pt,
                  font=\tiny\bfseries\color{UCRmgray}, minimum height=0.55cm}
  ]
    \node[box=UCRnavy] (interp) {\S3\ Interpretation\\[-1pt]for Safety};
    \node[box=UCRnavy, right=2.5cm of interp] (enhance) {\S4\ Safety\\[-1pt]Enhancements};
    \node[box=UCRnavy, right=2.5cm of enhance] (tools) {\S5\ Practical\\[-1pt]Tools};

    \draw[arr=UCRnavy] (interp.east) -- node[above,font=\scriptsize,text=UCRmgray]{informs} (enhance.west);
    \draw[arr=UCRnavy] (enhance.east) -- node[above,font=\scriptsize,text=UCRmgray]{operationalized by} (tools.west);

    \node[lbl, below=0.35cm of interp]
      {Training Attribution\\Input Tokens\\Model Internals\\Self-Reasoning};
    \node[lbl, below=0.35cm of enhance]
      {Steer Latent Vectors\\Modulate Neurons\\Edit Components\\Verify Outputs};
    \node[lbl, below=0.35cm of tools]
      {TDA Visualizers\\Token Vis\\Latent Vec Vis\\Neuron Explorers};

    \node[wfbox, above=0.6cm of enhance, minimum width=10.8cm]
      {LLM Workflow Lens:\ \ \ Training\ \ $\to$\ \ Input Tokens\ \ $\to$\ \ Model Inference\ \ $\to$\ \ Generation};
  \end{tikzpicture}

  \vspace{0.5em}
  \begin{columns}[T]
    \column{0.48\textwidth}
    {\small\bfseries\navy{Four Safety Concerns Covered}}
    \begin{itemize}\scriptsize\setlength\itemsep{3pt}
      \item \textbf{Hallucination}: confident fabrication of false facts
      \item \textbf{Jailbreaks \& Harmfulness}: bypassing safety alignment
      \item \textbf{Bias}: skewed or discriminatory outputs
      \item \textbf{Privacy Leakage}: memorized sensitive training data
    \end{itemize}
    \column{0.48\textwidth}
    \begin{ucrinfo}
      \scriptsize\bfseries Why these four share a root cause:\\[3pt]
      \normalfont\scriptsize
      All stem from training data with risky patterns + objectives that favour memorization over generalization. Interpretation reveals \emph{where} these manifest in model internals.
    \end{ucrinfo}
  \end{columns}
\end{frame}


% ─────────────────────────────────────────────────────────────────────────────
%  SECTION I
% ─────────────────────────────────────────────────────────────────────────────
\sectionframe{3}{Interpretation Methods for Safety}{How Do We Look Inside an LLM?}{Training Attribution · Input Tokens · Model Internals · Self-Reasoning}


% ═══════════════════════════════════════════════════════════════════
%  SLIDE 3 --- §3.1 TRAINING DATA ATTRIBUTION
% ═══════════════════════════════════════════════════════════════════
\begin{frame}{\S3.1\ Training Data Attribution (TDA)}
  {\footnotesize\itshape Core hypothesis: unsafe outputs stem from problematic training data.\\
  Can we trace which examples caused them?}
  \vspace{0.3em}
  \begin{columns}[T]
    \column{0.48\textwidth}
    {\small\bfseries\navy{Three Attribution Approaches}}
    \begin{ucrcard}[UCRnavy]{Representation-Based}
      \scriptsize Compare similarity of training example latent vectors to the output. Fast, but does \textbf{not establish causality}.
    \end{ucrcard}\vspace{2pt}
    \begin{ucrcard}[UCRnavy]{Gradient-Based (TracIn / Influence Functions)}
      \scriptsize Estimates how removing a training point would change model parameters. Theoretically grounded --- but \textbf{expensive to scale} to LLMs and sensitive to convexity assumptions.
    \end{ucrcard}\vspace{2pt}
    \begin{ucrcard}[UCRnavy]{Data Shapley}
      \scriptsize Marginal-contribution based attribution. Principled, but \textbf{limited to smaller models} due to computational cost.
    \end{ucrcard}

    \column{0.48\textwidth}
    {\small\bfseries\navy{Safety Applications}}
    \begin{itemize}\scriptsize\setlength\itemsep{4pt}
      \item \textbf{Hallucination sources}: \textit{LLM Attributor} (Lee et al., 2025) traces confabulated facts directly back to training examples --- shown as an interactive visualization
      \item \textbf{Bias origins}: identifies which web-scraped documents inject gender or racial stereotypes
      \item \textbf{Toxic completions}: links toxic outputs to specific training passages
      \item \textbf{Training dynamics}: tracks how safety capabilities like ``toxic content refusal'' emerge across training checkpoints
    \end{itemize}
    \vspace{0.3em}
    \begin{ucrwarn}
      \scriptsize\bfseries Open Challenge (\S6.3):\\[2pt]
      \normalfont\scriptsize
      Proprietary training data is inaccessible. Using TDA to actively \emph{enhance} safety (remove bad data + retrain) has only been shown for small, non-generative models. Scaling is unsolved.
    \end{ucrwarn}
  \end{columns}
\end{frame}


% ═══════════════════════════════════════════════════════════════════
%  SLIDE 4 --- §3.2 INPUT TOKEN IDENTIFICATION
% ═══════════════════════════════════════════════════════════════════
\begin{frame}{\S3.2\ Identifying Safety-Critical Input Tokens}
  {\footnotesize\itshape Which tokens in the prompt are responsible for triggering unsafe outputs?}
  \vspace{0.3em}
  \begin{columns}[T]
    \column{0.48\textwidth}
    {\small\bfseries\navy{Attribution Method Families}}
    \begin{itemize}\scriptsize\setlength\itemsep{4pt}
      \item \textbf{Attention Weights}: early and intuitive --- higher weight $\approx$ more important. Aggregation (rollout, mean/max) reduces overhead across heads. \muted{But: attention $\ne$ importance.}
      \item \textbf{Gradient-Based}: compute $\partial \text{output}/\partial \text{input embedding}$ --- how much does each token shift the output? More faithful than raw attention.
      \item \textbf{Perturbation-Based}: mask/replace tokens and observe output changes. Model-agnostic, but creates out-of-distribution inputs.
      \item \textbf{Shapley Values}: fair attribution across token subsets. Extended to LLMs via TokenSHAP, TextGenSHAP.
    \end{itemize}

    \column{0.48\textwidth}
    {\small\bfseries\navy{Safety Applications}}
    \begin{ucrcard}[UCRnavy]{Jailbreak Token Detection}
      \scriptsize
      \textbf{Cohen-Wang et al.\ (2024)}: perturbation-based attribution identifies the exact tokens triggering prompt poisoning attacks.\\[3pt]
      \textbf{Wang et al.\ (2025)}: prompt LLMs to self-identify which input tokens activate jailbreak circuits --- then remove them before generation.
    \end{ucrcard}
    \vspace{0.3em}
    \begin{ucrcard}[UCRnavy]{Bias Diagnosis}
      \scriptsize
      \textbf{Mohammadi (2024)}: Shapley values identify which tokens introduce gendered or racial bias --- enabling targeted suppression without changing non-biased outputs.
    \end{ucrcard}
    \vspace{0.3em}
    \begin{ucrdark}[UCRnavy]
      \centering\scriptsize
      Token attribution turns safety fixes into a \textbf{surgical filter}: remove only the offending tokens, preserve the rest of the user's intent.
    \end{ucrdark}
  \end{columns}
\end{frame}


% ═══════════════════════════════════════════════════════════════════
%  SLIDE 5 --- §3.3.1 PROBING LATENT SPACE
% ═══════════════════════════════════════════════════════════════════
\begin{frame}{\S3.3.1\ Probing Safety Regions in Latent Space}
  {\footnotesize\itshape \textbf{Linear Representation Hypothesis}: safety-relevant concepts (harmfulness, truthfulness, bias) are encoded as linear directions in the model's latent space.}
  \vspace{0.3em}
  \begin{columns}[T]
    \column{0.48\textwidth}
    {\small\bfseries\navy{Probing Techniques}}
    \begin{ucrcard}[UCRnavy]{Mean Difference Vectors}
      \scriptsize
      Compute mean latent vector for harmful vs.\ harmless prompts. Their difference encodes the ``harmfulness direction.''\\[2pt]
      \muted{Zou et al.\ (2023): Representation Engineering}
    \end{ucrcard}\vspace{2pt}
    \begin{ucrcard}[UCRnavy]{Probing Classifiers}
      \scriptsize
      Train a linear classifier on latent vectors to predict safety properties. Successfully detects hallucinations, jailbreaks, and bias.\\[2pt]
      \muted{Known issue: poor cross-task generalization.}
    \end{ucrcard}\vspace{2pt}
    \begin{ucrcard}[UCRnavy]{PCA / SVD on Activations}
      \scriptsize
      Dimensionality reduction uncovers the principal axes of unsafe behavior (Ball et al., 2024; Duan et al., 2024).
    \end{ucrcard}

    \column{0.48\textwidth}
    {\small\bfseries\navy{Key Safety Discoveries}}
    \begin{ucrcard}[UCRnavy]{The Refusal Direction (Arditi et al., 2024)}
      \scriptsize
      Jailbreak resistance is mediated by a \textbf{single linear direction} in activation space. Ablating it causes the model to comply with almost any harmful request.\\[3pt]
      \textbf{Implication}: robust safety = preserving this direction under adversarial pressure.
    \end{ucrcard}
    \vspace{0.3em}
    \begin{ucrcard}[UCRnavy]{Universal Truthfulness Hyperplane (Liu et al., 2024)}
      \scriptsize
      A separating hyperplane exists in LLM latent space that distinguishes truthful from hallucinated generations --- identifiable \emph{before} the output is produced.\\[3pt]
      Enables real-time hallucination interception.
    \end{ucrcard}
    \vspace{0.2em}
    \begin{ucrdark}[UCRnavy]
      \centering\scriptsize
      Probing is the \textbf{most widely used} interpretation method for safety --- enabling real-time detection of unsafe generation in a single forward pass.
    \end{ucrdark}
  \end{columns}
\end{frame}


% ═══════════════════════════════════════════════════════════════════
%  SLIDE 6 --- §3.3.2 PERTURBING COMPONENTS
% ═══════════════════════════════════════════════════════════════════
\begin{frame}{\S3.3.2\ Perturbing Components to Find Safety Mechanisms}
  {\footnotesize\itshape Knock out model components and observe what breaks --- establishing causal attribution of unsafe behavior.}
  \vspace{0.3em}
  \begin{columns}[T]
    \column{0.48\textwidth}
    {\small\bfseries\navy{Perturbation Methods}}
    \begin{itemize}\scriptsize\setlength\itemsep{4pt}
      \item \textbf{Component Knockout}: ablate layers, attention heads, or parameters to identify which are responsible for hallucinations, jailbreaks, and bias.
      \item \textbf{Activation Patching}: replace intermediate activations from one input with another input's activations --- directly measures causal effect of a specific region.
      \item \textbf{Circuit Extraction}: build graphs of important components (nodes) and information flow (edges). Automates mechanistic analysis.
      \item \textbf{Parameter Perturbation}: modify weights; evaluate how safety properties change --- tests alignment robustness.
    \end{itemize}
    \vspace{0.2em}
    \begin{ucrwarn}
      \scriptsize\bfseries Current Limitation:\\[2pt]
      \normalfont\scriptsize
      Most circuit analysis targets simple grammar or arithmetic tasks. Very few studies address real-world safety problems at scale (Hanna et al., 2024).
    \end{ucrwarn}

    \column{0.48\textwidth}
    {\small\bfseries\navy{Safety-Relevant Findings}}
    \begin{ucrcard}[UCRnavy]{Jailbreak-Critical Attention Heads}
      \scriptsize
      \textbf{Zhao et al.\ (2024)} and \textbf{Wei et al.\ (2024)}: specific attention heads mediate safety alignment. Pruning them enables jailbreaks; strengthening them improves robustness. Safety can be localized to specific heads per layer.
    \end{ucrcard}
    \vspace{0.3em}
    \begin{ucrcard}[UCRnavy]{Bias Localization}
      \scriptsize
      \textbf{Yang et al.\ (2023)} and \textbf{Ma et al.\ (2023)}: component knockout pinpoints exact attention heads and MLP sublayers encoding gender and racial bias --- enabling targeted debiasing without degrading other capabilities.
    \end{ucrcard}
    \vspace{0.3em}
    \begin{ucrcard}[UCRnavy]{RAG Hallucination Circuits}
      \scriptsize
      \textbf{Jin et al.\ (2024)}: activation patching localizes components responsible for ignoring retrieved context --- explaining RAG faithfulness failures.
    \end{ucrcard}
  \end{columns}
\end{frame}


% ═══════════════════════════════════════════════════════════════════
%  SLIDE 7 --- §3.3.3 SAEs + LOGIT LENS
% ═══════════════════════════════════════════════════════════════════
\begin{frame}{\S3.3.3\ Sparse Autoencoders: Deciphering What Neurons Encode}
  {\footnotesize\itshape Most neurons are \textbf{polysemantic} --- they encode multiple unrelated concepts simultaneously. SAEs disentangle them into interpretable features.}
  \vspace{0.3em}
  \begin{columns}[T]
    \column{0.48\textwidth}
    {\small\bfseries\navy{How SAEs Work}}
    \begin{enumerate}\scriptsize\setlength\itemsep{3pt}
      \item Encoder maps a polysemantic activation $\to$ high-dimensional \textbf{sparse} concept vector
      \item Each dimension (SAE ``neuron'') = one distinct, interpretable concept
      \item Decoder reconstructs the original activation
      \item Researchers find inputs that maximally activate each SAE neuron $\to$ human-readable labels
    \end{enumerate}
    \vspace{0.3em}
    \begin{ucrdark}[UCRnavy]
      \scriptsize\bfseries Anthropic's Scaling Monosemanticity (2024):\\[3pt]
      \normalfont\scriptsize
      SAEs applied to Claude 3 Sonnet. Millions of interpretable features found, including: deception, sycophancy, bias, dangerous content, and the famous \emph{``Golden Gate Bridge'' neuron} covered by your teammates.
    \end{ucrdark}

    \column{0.48\textwidth}
    {\small\bfseries\navy{Safety-Relevant SAE Findings}}
    \begin{ucrcard}[UCRnavy]{Jailbreak Features}
      \scriptsize
      \textbf{Härle et al.\ (2024)}: SCAR finds SAE neurons activating during jailbreak prompts. Suppressing them blocks the attack \textbf{without any retraining}.
    \end{ucrcard}
    \vspace{0.2em}
    \begin{ucrcard}[UCRnavy]{Bias Features}
      \scriptsize
      \textbf{Hegde (2024)} and \textbf{Zhou et al.\ (2025)}: specific SAE neurons encode gender bias. Clamping them to zero substantially reduces biased outputs.
    \end{ucrcard}
    \vspace{0.2em}
    \begin{ucrcard}[UCRnavy]{Privacy Leakage Features}
      \scriptsize
      \textbf{Frikha et al.\ (2025)}: PrivacyScalpel locates neurons encoding memorized personal information --- surgical privacy protection.
    \end{ucrcard}
    \vspace{0.2em}
    \begin{ucrinfo}
      \scriptsize\bfseries Also: Logit Lens\\[2pt]
      \normalfont\scriptsize
      Projects intermediate hidden states onto the vocabulary --- reveals what the model is ``thinking'' at each layer. Used to catch hallucinations forming mid-generation before the final token is produced.
    \end{ucrinfo}
  \end{columns}
\end{frame}


% ═══════════════════════════════════════════════════════════════════
%  SLIDE 8 --- §3.4 SELF-REASONING
% ═══════════════════════════════════════════════════════════════════
\begin{frame}{\S3.4\ Self-Reasoning: LLMs Interpreting Their Own Outputs}
  {\footnotesize\itshape LLMs can be prompted or trained to explain their own behavior --- but reliability is not guaranteed.}
  \vspace{0.3em}
  \begin{columns}[T]
    \column{0.48\textwidth}
    {\small\bfseries\navy{In-Generation Reasoning}}
    \begin{ucrcard}[UCRnavy]{Chain-of-Thought (CoT)}
      \scriptsize
      Prompting LLMs to generate intermediate reasoning steps before answering.\\[4pt]
      \textbf{Safety Use}: CoT surfaces potential harms before the model commits to a response --- a form of runtime self-audit. Many variants: Tree-of-Thought, Graph-of-Thought, Faithful CoT.
    \end{ucrcard}
    \vspace{0.3em}
    \begin{ucrcard}[UCRnavy]{Uncertainty Estimation}
      \scriptsize
      Models trained or prompted to express confidence levels. Flags when hallucination is likely --- enables early warning systems for downstream users.
    \end{ucrcard}

    \column{0.48\textwidth}
    {\small\bfseries\navy{Post-Hoc Self-Explanation}}
    \begin{ucrcard}[UCRnavy]{Chain-of-Verification (Dhuliawala et al., 2024)}
      \scriptsize
      Split a response into factual claims, then ask the model to verify each against its own knowledge. Iterative self-questioning substantially reduces hallucination rates on tasks like WikiData QA.
    \end{ucrcard}
    \vspace{0.3em}
    \begin{ucrwarn}
      \scriptsize\bfseries Critical Caveat: Unfaithful Explanations\\[3pt]
      \normalfont\scriptsize
      CoT explanations often \textbf{do not reflect} the model's actual computation (Turpin et al., 2023). Models confidently produce plausible-but-wrong self-explanations.\\[3pt]
      $\therefore$ Self-reasoning is a \textbf{useful signal}, not a guarantee --- always needs external verification.
    \end{ucrwarn}
  \end{columns}
\end{frame}


% ─────────────────────────────────────────────────────────────────────────────
%  SECTION II
% ─────────────────────────────────────────────────────────────────────────────
\sectionframe{4}{Enhancing Safety via Interpretation}{From Understanding to Action}{Steer · Modulate · Edit · Verify · Self-Reason}


% ═══════════════════════════════════════════════════════════════════
%  SLIDE 9 --- §4.1 + §4.2.1 STEERING LATENT VECTORS
% ═══════════════════════════════════════════════════════════════════
\begin{frame}{\S4.1 \& \S4.2.1\ Token-Level Fixes and Representation Steering}
  \begin{columns}[T]
    \column{0.48\textwidth}
    {\small\bfseries\navy{\S4.1 --- Attend to Relevant Tokens}}
    \begin{itemize}\scriptsize\setlength\itemsep{3pt}
      \item Remove jailbreak-triggering tokens identified by attribution (\textit{Pan et al., 2025})
      \item Amplify attention to reliable user-specified tokens (\textit{Zhang et al., 2024})
      \item Redirect model focus to relevant context to reduce RAG hallucinations (\textit{Liu et al., 2025})
    \end{itemize}
    \vspace{0.5em}
    {\small\bfseries\navy{\S4.2.1 --- Representation Engineering}}
    \begin{ucrcard}[UCRnavy]{Steering Vectors}
      \scriptsize
      Add the ``safe direction'' (found by probing) to the model's hidden states at inference time:\\[4pt]
      $\qquad h_{\text{new}} = h + \alpha \cdot \hat{v}_{\text{safe}}$\\[4pt]
      \textbf{No retraining.} Works at inference time. Reversible. Can be tuned per task.
    \end{ucrcard}

    \column{0.48\textwidth}
    {\small\bfseries\navy{Key Results}}
    \begin{ucrcard}[UCRnavy]{Inference-Time Intervention (ITI) --- Li et al., 2023}
      \scriptsize
      Probes identify truthfulness directions in Llama activations. Amplifying them at inference time achieves $+16\%$ on TruthfulQA with \textbf{no fine-tuning} --- one of the first demonstrations of mechanistic safety improvement.
    \end{ucrcard}
    \vspace{0.3em}
    \begin{ucrcard}[UCRnavy]{Contrastive Activation Addition --- Rimsky et al., 2024}
      \scriptsize
      Subtract the ``sycophancy'' or ``bias'' direction from hidden states of Llama 2. Effective and simple --- requires \textbf{zero fine-tuning}, works across many behaviors.
    \end{ucrcard}
    \vspace{0.3em}
    \begin{ucrdark}[UCRnavy]
      \centering\scriptsize
      Latent steering: parameter-efficient, fast, directly grounded in mechanistic understanding of \emph{where} the unsafe behavior lives.
    \end{ucrdark}
  \end{columns}
\end{frame}


% ═══════════════════════════════════════════════════════════════════
%  SLIDE 10 --- §4.2.2 MODULATING NEURONS (SAE-BASED)
% ═══════════════════════════════════════════════════════════════════
\begin{frame}{\S4.2.2\ Modulating Neurons: SAE-Based Safety Control}
  {\footnotesize\itshape Suppress risky SAE neurons or amplify safe ones to guide the model away from unsafe behaviors --- without any weight updates.}
  \vspace{0.3em}
  \begin{columns}[T]
    \column{0.48\textwidth}
    {\small\bfseries\navy{The Mechanism}}
    \begin{enumerate}\scriptsize\setlength\itemsep{4pt}
      \item Run the SAE encoder to get a sparse concept vector for the current activation
      \item Identify the SAE neuron(s) corresponding to the target concept (e.g., ``deception'', ``CSAM'', ``violence'')
      \item \textbf{Clamp, suppress, or amplify} those neurons in the sparse vector
      \item Pass the modified concept vector through the SAE decoder to reconstruct a ``safer'' activation
      \item Continue generation from this modified hidden state
    \end{enumerate}
    \vspace{0.3em}
    \begin{ucrdark}[UCRnavy]
      \centering\scriptsize\bfseries
      No fine-tuning. No full retraining.\\
      Targeted · Reversible · Interpretable.
    \end{ucrdark}

    \column{0.48\textwidth}
    {\small\bfseries\navy{Applications \& Results}}
    \begin{ucrcard}[UCRnavy]{Refusal Steering --- O'Brien et al., 2024}
      \scriptsize
      Identified the SAE neuron controlling \textbf{refusal behavior}. Amplifying it increases refusals on harmful requests. Suppressing it achieves near-universal compliance --- a clear dual-use risk showing how powerful this control is.
    \end{ucrcard}
    \vspace{0.3em}
    \begin{ucrcard}[UCRnavy]{Hallucination Mitigation --- SAFE (Abdaljalil et al., 2025)}
      \scriptsize
      SAE-based query enrichment suppresses hallucination-prone neurons in a RAG pipeline --- improves factual accuracy without knowledge base changes.
    \end{ucrcard}
    \vspace{0.3em}
    \begin{ucrwarn}
      \scriptsize\bfseries Dual-Use Warning:\\[2pt]
      \normalfont\scriptsize
      The same SAE control mechanisms that enhance safety can amplify harmful features or suppress safety features. Requires careful access controls and auditing.
    \end{ucrwarn}
  \end{columns}
\end{frame}


% ═══════════════════════════════════════════════════════════════════
%  SLIDE 11 --- §4.2.3 EDITING HARMFUL COMPONENTS
% ═══════════════════════════════════════════════════════════════════
\begin{frame}{\S4.2.3\ Editing Harmful Model Components}
  {\footnotesize\itshape Directly modify the weights of safety-critical components --- pruning, direction editing, and knowledge surgery.}
  \vspace{0.3em}
  \begin{columns}[T]
    \column{0.48\textwidth}
    {\small\bfseries\navy{Approach Family}}
    \begin{ucrcard}[UCRnavy]{Pruning / Downscaling}
      \scriptsize Remove or downweight attention heads / sublayers identified as responsible for harmful outputs. Layer-specific jailbreak defense (Zhao et al., 2024), hallucination heads (Li et al., 2024), bias heads (Ma et al., 2023).
    \end{ucrcard}\vspace{2pt}
    \begin{ucrcard}[UCRnavy]{Safety Direction Editing}
      \scriptsize Compare aligned vs.\ unaligned model weights $\to$ find the ``safety direction'' in parameter space $\to$ steer critical parameters accordingly. \textit{Safety Arithmetic} (Hazra et al., 2024).
    \end{ucrcard}\vspace{2pt}
    \begin{ucrcard}[UCRnavy]{Knowledge Editing (ROME / MEMIT)}
      \scriptsize Edit specific factual associations in model weights without retraining. Can remove dangerous knowledge or correct biased associations at the fact level.
    \end{ucrcard}

    \column{0.48\textwidth}
    {\small\bfseries\navy{Key Results}}
    \begin{ucrcard}[UCRnavy]{Safety Layers (Li et al., ICLR 2025)}
      \scriptsize
      A small set of ``safety layers'' exists in aligned LLMs. \textbf{Protecting these layers from fine-tuning} preserves safety even when the model is further trained on task-specific data --- directly addresses the fine-tuning attack problem.
    \end{ucrcard}
    \vspace{0.3em}
    \begin{ucrcard}[UCRnavy]{Layer-Specific Jailbreak Defense (Zhao et al., EMNLP 2024)}
      \scriptsize
      Identifies and modifies the exact layers where jailbreaks take effect. Substantially improves robustness on StrongREJECT without degrading helpfulness on benign queries.
    \end{ucrcard}
    \vspace{0.3em}
    \begin{ucrdark}[UCRnavy]
      \centering\scriptsize\bfseries Tradeoff: Precision vs.\ Side-Effects\\[2pt]
      \normalfont\scriptsize
      Editing too aggressively degrades model utility. Finding the minimum effective edit with maximum safety gain is an active research frontier.
    \end{ucrdark}
  \end{columns}
\end{frame}


% ═══════════════════════════════════════════════════════════════════
%  SLIDE 12 --- §4.3 + §4.4 VERIFY & SELF-REASONING OUTPUT
% ═══════════════════════════════════════════════════════════════════
\begin{frame}{\S4.3 \& \S4.4\ Verify Before Output + Self-Reasoning Safety}
  \begin{columns}[T]
    \column{0.48\textwidth}
    {\small\bfseries\navy{\S4.3 --- Verify Before Output}}
    \begin{ucrcard}[UCRnavy]{Best-of-N with Safety Filtering}
      \scriptsize
      Generate multiple candidate responses. Use probing classifiers or self-evaluation to select only safe ones.\\[3pt]
      \textbf{SelfCheck} (Miao et al., 2024): LLMs zero-shot check their own reasoning for internal inconsistency.\\[3pt]
      \textbf{Chain-of-Verification} (Dhuliawala et al., 2024): iterative fact-checking reduces hallucination substantially.
    \end{ucrcard}
    \vspace{0.3em}
    \begin{ucrcard}[UCRnavy]{Intervention Mid-Generation}
      \scriptsize
      Detect unsafe token sequences using hidden-state probes \emph{during} generation and \textbf{resample} from a safer distribution before committing to the harmful text.
    \end{ucrcard}

    \column{0.48\textwidth}
    {\small\bfseries\navy{\S4.4 --- Output with Self-Reasoning}}
    \begin{ucrcard}[UCRnavy]{GuardReasoner (Liu et al., 2025)}
      \scriptsize
      Fine-tune LLMs to generate intermediate reasoning explaining \emph{why} a response may be harmful before deciding whether to comply. Makes safety decisions \textbf{transparent and auditable}.
    \end{ucrcard}
    \vspace{0.3em}
    \begin{ucrcard}[UCRnavy]{CoT-Based Bias \& Safety Reasoning}
      \scriptsize
      \textbf{Kaneko et al.\ (2024)}: CoT prompting for gender bias evaluation --- model articulates fairness reasoning explicitly.\\[3pt]
      \textbf{Cao et al.\ (2024)}: CoT prompting for jailbreak defense --- chain of thought acts as an in-context safety guardrail.
    \end{ucrcard}
    \vspace{0.2em}
    \begin{ucrdark}[UCRnavy]
      \centering\scriptsize
      Self-reasoning safety methods produce \textbf{auditable} refusals: the decision comes with an explanation, not just a silent block.
    \end{ucrdark}
  \end{columns}
\end{frame}


% ─────────────────────────────────────────────────────────────────────────────
%  SECTION III
% ─────────────────────────────────────────────────────────────────────────────
\sectionframe{5}{Practical Tools}{Operationalizing Safety-Focused Interpretation}{Visualization · Interaction · Actionability}


% ═══════════════════════════════════════════════════════════════════
%  SLIDE 13 --- §5 TOOLS
% ═══════════════════════════════════════════════════════════════════
\begin{frame}{\S5\ Tools: Making Interpretation Actionable}
  {\footnotesize\itshape Interpretation methods only matter if practitioners can act on them. Four tool categories bridge theory to practice.}
  \vspace{0.3em}
  \begin{columns}[T]
    \column{0.48\textwidth}
    \begin{ucrcard}[UCRnavy]{\S5.1\ Training Data Attribution Visualizers}
      \scriptsize
      \textbf{LLM Attributor} (Lee et al., 2025): interactive tool tracing model outputs to training data. Users inspect which examples caused a hallucinated fact or toxic output --- and can flag them for removal.
    \end{ucrcard}\vspace{3pt}
    \begin{ucrcard}[UCRnavy]{\S5.2\ Input Token Visualizations}
      \scriptsize
      Tools like \textbf{BertViz} (Vig, 2019), \textbf{iScore}, \textbf{PromptAID} (Mishra et al., 2025): visualize attention patterns across heads, reveal spurious token associations indicating bias. Support \textbf{interactive perturbation} --- edit tokens, observe output shift live.
    \end{ucrcard}

    \column{0.48\textwidth}
    \begin{ucrcard}[UCRnavy]{\S5.3\ Latent Vector Visualizations}
      \scriptsize
      Project latent vectors into 2D space; \textbf{logit lens} visualizers show vocabulary-level semantics at each layer; interactive steering tools let users push activations toward safer regions live. \textit{LM-Debugger} (Geva et al., 2022); \textit{Dashboard for Transparency} (Chen et al., 2024).
    \end{ucrcard}\vspace{3pt}
    \begin{ucrcard}[UCRnavy]{\S5.4\ Neuron \& SAE Visualizations}
      \scriptsize
      \textbf{Neuronpedia} (Lin, 2023): interactive reference for SAE neuron concepts. \textbf{Neuroscope} (Nanda, 2022): browser for mechanistic interpretability of LLMs. Enable safety researchers to find, inspect, and manipulate safety-relevant features at scale.
    \end{ucrcard}
    \vspace{0.2em}
    \begin{ucrinfo}
      \centering\scriptsize
      Common thread: all categories enable \textbf{interactive} exploration --- practitioners can act on findings, not just read about them.
    \end{ucrinfo}
  \end{columns}
\end{frame}


% ─────────────────────────────────────────────────────────────────────────────
%  SECTION IV --- OPEN CHALLENGES
% ─────────────────────────────────────────────────────────────────────────────
\sectionframe{6}{Open Challenges \& Future Directions}{What Remains Unsolved?}{Attacks · Evaluation · Training Attribution · Presentation · Safety Scope}


% ═══════════════════════════════════════════════════════════════════
%  SLIDE 14 --- OPEN CHALLENGES
% ═══════════════════════════════════════════════════════════════════
\begin{frame}{Open Challenges at the Interpretability–Safety Interface}
  \begin{columns}[T]
    \column{0.48\textwidth}
    \begin{ucrcard}[UCRnavy]{[!] Interpretation-Based Attacks}
      \scriptsize
      Interpretation reveals \emph{why} certain tokens trigger unsafe outputs --- enabling adversaries to craft \textbf{more targeted} jailbreak suffixes (Lin et al., 2024; Winninger et al., 2025).\\[3pt]
      Adversaries use SAE feature maps to identify and manipulate model components linked to safety alignment (Arditi et al., 2024).\\[3pt]
      \textbf{Tools built for transparency become attack vectors.}
    \end{ucrcard}\vspace{3pt}
    \begin{ucrcard}[UCRnavy]{[?] Reliable Evaluation of Interpretations}
      \scriptsize
      No field consensus on ``faithful'' or ``interpretable.'' Evaluation scores vary with \emph{presentation format}. Requires interdisciplinary collaboration: ML $+$ HCI $+$ cognitive science.
    \end{ucrcard}

    \column{0.48\textwidth}
    \begin{ucrcard}[UCRnavy]{[DB] TDA for Safety Enhancement}
      \scriptsize
      TDA can trace unsafe behavior to training data --- but \textbf{using that information to fix safety} (retrain without bad data) only demonstrated on small, non-generative models. Scaling to LLMs is an open problem.
    \end{ucrcard}\vspace{3pt}
    \begin{ucrcard}[UCRnavy]{[People] User-Centered Presentation}
      \scriptsize
      Conversational interaction (e.g., TalkToModel) aids human understanding --- but \textbf{no tools} apply this to safety-oriented interpretation. Different stakeholders (developers, auditors, regulators) need different views.
    \end{ucrcard}\vspace{3pt}
    \begin{ucrcard}[UCRnavy]{[Scope] Broader Safety Dimensions}
      \scriptsize
      Interpretation research focuses on hallucination, jailbreaks, bias, and privacy. \textbf{Under-studied}: code security, OOD robustness, and over-refusal (Siska \& Sankaran, 2025).
    \end{ucrcard}
  \end{columns}
\end{frame}


% ═══════════════════════════════════════════════════════════════════
%  SLIDE 15 --- FULL PIPELINE SYNTHESIS
% ═══════════════════════════════════════════════════════════════════
\begin{frame}[fragile]{Putting It All Together: The Full Interpret $\to$ Fix Pipeline}
  \vspace*{0.2cm}
  \begin{tikzpicture}[
    stage/.style={rectangle, draw=#1, fill=#1!12, text=#1,
                  font=\tiny\bfseries, minimum width=2.3cm,
                  minimum height=0.7cm, align=center, rounded corners=2pt},
    fix/.style={rectangle, draw=UCRnavy!60, fill=UCRnavy!8, font=\tiny,
                text width=2.2cm, align=center, rounded corners=1pt, minimum height=0.85cm,
                text=UCRdark},
    arr/.style={-{Stealth[length=4pt]}, thick, #1},
    lbl/.style={font=\tiny\itshape, text=UCRmgray},
  ]
    % Interpretation row
    \node[stage=UCRnavy] (tda)   {Training\\Attribution};
    \node[stage=UCRnavy, right=0.5cm of tda]   (token) {Input Token\\Attribution};
    \node[stage=UCRnavy, right=0.5cm of token] (probe) {Latent Space\\Probing};
    \node[stage=UCRnavy, right=0.5cm of probe]   (sae)   {SAE / Neuron\\Deciphering};
    \node[stage=UCRnavy, right=0.5cm of sae]  (self)  {Self-Reasoning\\Explanation};

    \draw[arr=UCRmgray] (tda) -- (token);
    \draw[arr=UCRmgray] (token) -- (probe);
    \draw[arr=UCRmgray] (probe) -- (sae);
    \draw[arr=UCRmgray] (sae) -- (self);

    % Fix row below
    \node[fix, below=0.55cm of tda]   (f1) {Remove / reweight\\bad training data};
    \node[fix, below=0.55cm of token] (f2) {Suppress jailbreak\\trigger tokens};
    \node[fix, below=0.55cm of probe] (f3) {Steer latent vectors\\to safe region};
    \node[fix, below=0.55cm of sae]   (f4) {Clamp unsafe\\SAE neurons};
    \node[fix, below=0.55cm of self]  (f5) {Verify output;\\reason about safety};

    \draw[arr=UCRnavy] (tda) -- (f1);
    \draw[arr=UCRnavy] (token) -- (f2);
    \draw[arr=UCRnavy] (probe) -- (f3);
    \draw[arr=UCRnavy] (sae) -- (f4);
    \draw[arr=UCRnavy] (self) -- (f5);

    % Section labels
    \node[lbl, below=0.1cm of f1] {\S4.2.3};
    \node[lbl, below=0.1cm of f2] {\S4.1};
    \node[lbl, below=0.1cm of f3] {\S4.2.1};
    \node[lbl, below=0.1cm of f4] {\S4.2.2};
    \node[lbl, below=0.1cm of f5] {\S4.3--4.4};
  \end{tikzpicture}

  \vspace{0.4em}
  \begin{columns}[T]
    \column{0.48\textwidth}
    {\small\bfseries\navy{What Makes This Approach Powerful}}
    \begin{itemize}\scriptsize\setlength\itemsep{2pt}
      \item Mechanistic understanding, not just empirical patches
      \item Many methods work \textbf{at inference time} --- no retraining
      \item Fixes are \textbf{targeted}: surgery, not chemotherapy
      \item Produces \textbf{auditable} safety decisions with explanations
    \end{itemize}

    \column{0.48\textwidth}
    {\small\bfseries\navy{What Are the Risks \& Limits}}
    \begin{itemize}\scriptsize\setlength\itemsep{2pt}
      \item Dual-use: interpretation tools can enhance \emph{or} undermine safety
      \item CoT explanations may not reflect actual model computation
      \item Scaling mechanistic circuit analysis to real tasks remains hard
      \item Non-linear safety representations challenge linear steering
    \end{itemize}
  \end{columns}
\end{frame}


% ═══════════════════════════════════════════════════════════════════
%  SLIDE 16 --- CONCLUSION
% ═══════════════════════════════════════════════════════════════════
\begin{frame}{Conclusion: Interpretation is the Path to Principled Safety}
  \vspace*{0.1cm}
  \begin{columns}[T]
    \column{0.48\textwidth}
    {\small\bfseries\navy{What We Covered}}
    \begin{enumerate}\scriptsize\setlength\itemsep{5pt}
      \item \textbf{\S3\ Interpretation Methods} (How we look inside)\\
        \muted{TDA · Token Attribution · Probing · Perturbation · SAEs · Self-Reasoning}
      \item \textbf{\S4\ Safety Enhancements} (How we fix it)\\
        \muted{Steering Vectors · Neuron Modulation · Model Editing · Output Verification}
      \item \textbf{\S5\ Practical Tools} (How we deploy it)\\
        \muted{TDA Vis · Token Vis · Latent Vis · Neuron/SAE Explorers}
      \item \textbf{\S6\ Open Challenges}\\
        \muted{Interpretation-based attacks · Eval · TDA scaling · Broader safety}
    \end{enumerate}

    \column{0.48\textwidth}
    \begin{ucrdark}[UCRnavy]
      \centering\small\bfseries
      The Core Message\\[6pt]
      \normalfont\scriptsize
      Unsafe LLM behavior is not random --- it has \textbf{identifiable, localizable causes} inside the model.\\[6pt]
      Interpretability gives us the tools to \textbf{find} those causes and \textbf{fix them precisely} --- transforming AI safety from an empirical art into a principled engineering discipline.
    \end{ucrdark}
    \vspace{0.4em}
    {\small\bfseries\navy{Key Takeaways}}
    \begin{itemize}\scriptsize\setlength\itemsep{3pt}
      \item Interpretability $\not=$ only explaining models; it \emph{enables} targeted safety fixes
      \item SAEs are becoming the workhorse of safety intervention --- but carry dual-use risk
      \item Much unexplored: TDA at scale, circuit analysis for safety, holistic safety ontologies
      \item The field needs the dual expertise of \textbf{both} your teammates' topics, working together
    \end{itemize}
  \end{columns}
\end{frame}


% ═══════════════════════════════════════════════════════════════════
%  BACKUP --- REFERENCE TABLE
% ═══════════════════════════════════════════════════════════════════
\begin{frame}{Reference: Interpretation × Safety Enhancement Summary}
  \vspace*{0.1cm}
  \small
  \setlength{\tabcolsep}{4pt}
  \renewcommand{\arraystretch}{1.45}
  \begin{tabularx}{\textwidth}{>{\bfseries\scriptsize\color{UCRnavy}}l >{\scriptsize}X >{\scriptsize}X}
    \rowcolor{UCRnavy}
    \color{white}Interpretation Method &
    \color{white}What It Reveals &
    \color{white}How It Enhances Safety \\
    \rowcolor{UCRlgray}
    Training Attribution & Poisonous training examples & Remove / reweight bad data (\S4.2.3) \\
    Input Token Attr. & Jailbreak trigger tokens & Suppress before generation (\S4.1) \\
    \rowcolor{UCRlgray}
    Latent Probing & Safety directions in activation space & Steer vectors; real-time detection (\S4.2.1) \\
    Component Perturbation & Harmful attention heads / layers & Prune / downscale components (\S4.2.3) \\
    \rowcolor{UCRlgray}
    SAEs & Polysemantic $\to$ monosemantic features & Clamp unsafe neurons precisely (\S4.2.2) \\
    Logit Lens & Layer-wise vocabulary predictions & Catch hallucinations mid-generation (\S4.3) \\
    \rowcolor{UCRlgray}
    Self-Reasoning (CoT) & Model's own safety assessment & Verify + explain before output (\S4.4) \\
  \end{tabularx}
  \vspace{0.4em}
  \begin{ucrinfo}
    \centering\scriptsize\bfseries
    Source: Lee, Cho, Kim, Peng, Phute, Chau (Georgia Tech) ·
    \textit{Interpretation Meets Safety} · EMNLP 2025
  \end{ucrinfo}
\end{frame}

